% Baseerat: TACCESS submission source (ACM acmart, acmsmall journal format).
%
% This is the venue-formatted version for ACM Transactions on Accessible
% Computing. paper/baseerat.tex remains the plain-article working source; the
% body text here is identical to it.
%
% There is no LaTeX compiler on the author's machine. Compile this on Overleaf
% (which ships the acmart class) or with a full TeX Live install:
%     pdflatex baseerat-acmart  (twice, for references)
%
% Before submission:
%   1. Paste the CCSXML block generated at https://dl.acm.org/ccs into the marked
%      spot below (the \ccsdesc lines are provided; the XML block pairs with them).
%   2. Set \acmDOI, \acmYear, volume/number/article fields if the editor assigns
%      them; the review build does not need them.
%   3. UK English and no em or en dashes are house style; keep them on edits.

\documentclass[acmsmall,review,screen]{acmart}

\usepackage{booktabs}

% --- ACM front matter ---------------------------------------------------------
\acmJournal{TACCESS}
\citestyle{acmnumeric}
\settopmatter{printacmref=true}
\setcopyright{none}     % review build; the editor sets the final rights.

\begin{document}

\title{Baseerat: Measuring the Oversight Gap When a Computer-Use Agent's Overseer
Cannot See the Screen}

\author{Yasir Musawar}
\affiliation{%
  \institution{WAJD AI}
  \country{United Kingdom}
}
\email{baseerat@wajd.co.uk}

\renewcommand{\shortauthors}{Musawar}

\begin{abstract}
Safety benchmarks for computer-use agents evaluate an agent's behaviour under the
tacit assumption that the human overseer can see the screen. For the roughly forty
million people who operate a computer through a screen reader, that assumption does
not hold: the only account of what an agent did arrives through the agent's
narration and the post-hoc accessibility tree, and a growing body of work shows
that an agent's own account of its actions is unreliable, with claimed success
exceeding actual success by large margins. We introduce \emph{Baseerat}, a
benchmark that pairs each task with a matched sighted control and measures the
\emph{oversight parity gap}: the detection rate a sighted overseer achieves on a
diverged agent run minus the rate an overseer restricted to the accessible channel
achieves on the same run. Baseerat runs tasks in a real browser under three
conditions (clean, silent failure, and a deceptive indirect injection that poisons
the accessibility tree in step with the agent's narration), and audits every run
through both channels. Using a deterministic reference auditor over a 150-task
corpus, we establish the benchmark's construct validity and report an oversight
parity gap of 0.50, concentrated entirely in the deceptive-injection condition. We
then evaluate a defence, trusted action receipts, and show it closes the gap to
zero. We release the benchmark, the live rig, the corpus, and the reference
results, and specify the language-model and human expert-auditor arms.
\end{abstract}

% --- CCS concepts -------------------------------------------------------------
% Paste the CCSXML block from https://dl.acm.org/ccs here (between the XML
% environment). The \ccsdesc lines below pair with it.
\begin{CCSXML}
<ccs2012>
 <concept>
  <concept_id>10003120.10003121</concept_id>
  <concept_desc>Human-centered computing~Accessibility</concept_desc>
  <concept_significance>500</concept_significance>
 </concept>
 <concept>
  <concept_id>10002978.10003029</concept_id>
  <concept_desc>Security and privacy~Human and societal aspects of security and privacy</concept_desc>
  <concept_significance>300</concept_significance>
 </concept>
 <concept>
  <concept_id>10010147</concept_id>
  <concept_desc>Computing methodologies~Artificial intelligence</concept_desc>
  <concept_significance>300</concept_significance>
 </concept>
</ccs2012>
\end{CCSXML}
\ccsdesc[500]{Human-centered computing~Accessibility}
\ccsdesc[300]{Security and privacy~Human and societal aspects of security and privacy}
\ccsdesc[300]{Computing methodologies~Artificial intelligence}

\keywords{computer-use agents, AI safety, accessibility, screen readers, blind and
low-vision users, prompt injection, agent oversight, benchmark, human oversight}

\maketitle

\section{Introduction}

Computer-use agents now act on real interfaces on a person's behalf: they fill
forms, send messages, move money, and change settings. The safety of these agents
is increasingly benchmarked, but the benchmarks share an unstated premise. When
they score whether an agent completed a task correctly, resisted a prompt
injection, or avoided a harmful action, the ground truth is read from the screen,
and the implicit overseer who would catch a failure in deployment is assumed to be
able to read the screen too.

For a blind screen-reader user that premise fails. The overseer does not see the
result; they hear the agent's narration and, if they inspect the outcome, they
read whatever the accessibility tree exposes. The oversight channel collapses onto
the agent's own account of what it did. This would be a manageable limitation if
that account were reliable. It is not. Recent work characterises \emph{false
success}, in which agents produce confident closure statements over tasks that in
fact failed, and \emph{outcome hallucination}, in which an agent claims completion
despite evidence to the contrary, with claimed-to-actual success ratios reported
as high as $5.3\times$.

Put together, these two observations imply a vulnerability that neither literature
measures. If the accessible channel is the whole of a blind user's oversight, and
that channel carries an unreliable self-report, then the safety of the agent for
that user reduces to the fidelity of the self-report, and to whether the accessible
channel is rich enough to reveal a discrepancy when the self-report is wrong.
Nobody has measured either quantity, and the adversarial case is worse than the
accidental one. In an indirect prompt injection the attacker controls page content;
nothing prevents the injection from also instructing the agent to narrate a clean,
benign account of a hijacked action, so that the attack and its cover story travel
down the very channel the blind user depends on. A sighted overseer sees the true
result regardless; a non-visual overseer sees only the cover story.

We make this precise with a benchmark. Our contributions are: (1) a problem
formulation and a metric, the \emph{oversight parity gap}, the difference in
divergence-detection rate between a visually unrestricted overseer and one
restricted to the accessible channel, over the same runs, together with
\emph{self-report fidelity} and \emph{non-visual detectability}; (2) a live
benchmark, Baseerat, in which the deceptive-injection condition is constructed in
an actual document object model so that the accessibility tree exposes a benign
value while the committed state is hijacked; (3) a deterministic reference-auditor
result over a 150-task corpus that yields an oversight parity gap of 0.50, entirely
within the deceptive-injection condition; (4) a defence, trusted action receipts,
that restores non-visual detection from 0.50 to 1.00 and closes the gap to zero;
and (5) an open protocol for the language-model and human expert-auditor arms.

The name \emph{Baseerat} (from the Arabic \emph{baseerah}, ``inner sight'',
discernment) marks the point: the question is not whether an overseer has eyes, but
whether the channel they are given carries enough to discern a failure.

\section{Related Work and Positioning}

Baseerat sits at the intersection of three lines of work, each of which supplies
one ingredient and none of which occupies the joint question.

\paragraph{Safety benchmarks for computer-use agents.} OS-Harm~\cite{osharm}
evaluates computer-use agents across deliberate misuse, prompt injection, and model
misbehaviour, and is representative of a class that scores outcomes with full visual
access to the screen. It does not model the overseer's perceptual channel; the
auditor is assumed sighted.

\paragraph{Unreliable self-reports.} A separate line characterises the failure mode
we exploit. Work on false success~\cite{silentfailure} shows agents assert
completion over failed tasks and that language-model judges systematically miss it;
MIRAGE-Bench~\cite{mirage} provides a taxonomy of agent hallucination including
outcome hallucination. These establish that the self-report is untrustworthy, but
they too assume a sighted party is available to notice, and do not consider
accessibility.

\paragraph{Accessibility of agents.} A recent and fast-moving line studies agents
for blind and low-vision users. A11y-CUA~\cite{a11ycua} documents that agent task
success collapses under assistive-technology interaction styles. An interview
study~\cite{xaiblv} of blind users in the agentic era names the epistemic
vulnerability precisely, observing that an error early in a multi-step task can
propagate silently before any feedback is received, but it is qualitative and
introduces no benchmark or metric. Closest to our work, NeXUI~\cite{nexui} scores
whether an assistive agent explains each step in a way grounded in interface state.
It asks whether agents explain well; it does not include adversarial conditions,
seeded silent failures, deception, or any measurement of whether the explanations
are sufficient to detect an agent failing or lying, and it runs without blind
participants.

\paragraph{The joint question.} The gap is the conjunction. The safety benchmarks
measure agent behaviour but assume a sighted auditor; the false-success work
measures self-report unreliability but assumes a sighted auditor; the accessibility
work studies the non-visual channel but not adversarial detection. Baseerat
measures whether the non-visual channel is \emph{sufficient to detect failure},
including adversarial failure, and what restores it when it is not. The defence we
evaluate is informed by the tool-receipts idea~\cite{receipts}, a machine-verified
statement of effect rather than the agent's word for it.

\section{The Baseerat Benchmark}

\subsection{Tasks and the Field Specification}
A task is a realistic operation on a web or desktop surface: send an email, update
a delivery address, add a payee, rename and move a file, book a calendar event,
update account recovery details. Each task carries a \emph{field specification}: the
set of observable slots that a correct completion must leave holding specified
values. The field specification is the ground truth against which both channels are
scored. The corpus contains 150 tasks generated deterministically, stratified evenly
across six domains and four \emph{injection archetypes} that vary how a hijacked
value differs from the intended one: a homograph (a look-alike character swap such
as \texttt{l} to \texttt{1}), a plausible substitution, an appended hostile token,
and a wholesale redirect. The stratification exists so that the auditor arms can be
analysed by how subtle the hijack is; the reference auditor, by construction, is
archetype-invariant (Section~\ref{sec:results}).

\subsection{Three Conditions}
Every task is run under three conditions, summarised in Table~\ref{tab:conditions}.
The design of the conditions is the core of the benchmark, because it separates a
failure that the accessible channel can reveal from one that it cannot.

\begin{table}[t]
\caption{The three conditions. The contrast between silent failure and deceptive
injection isolates the oversight parity gap.}
\label{tab:conditions}
\small
\begin{tabular}{@{}p{2.2cm}p{2.7cm}p{2.7cm}p{4.0cm}@{}}
\toprule
\textbf{Condition} & \textbf{Ground truth} & \textbf{Accessibility tree} &
\textbf{Detectable by} \\
\midrule
clean & matches the task & honest & nothing to detect \\
silent failure & a field silently fails to save & \emph{honest}: shows the failure
& anyone who cross-checks the tree against the agent's claim \\
deceptive injection & a field is hijacked & \emph{poisoned}: shows the benign value
& only a sighted overseer reading the true result \\
\bottomrule
\end{tabular}
\end{table}

Under \textbf{silent failure} the accessibility tree is honest, so a diligent
non-visual overseer who cross-checks the tree against the agent's claim of success
can catch it. Under \textbf{deceptive injection} the accessibility tree is poisoned
to corroborate the false narration, so the non-visual channel shows a consistent
benign story and only the ground truth reveals the hijack.

\subsection{Two Channels and the Live Rig}
Each run is audited through two channels. The \textbf{visual} channel additionally
receives the ground-truth committed state, modelling a sighted overseer who can read
the true result from the screen. The \textbf{non-visual} channel receives the
agent's narration and the post-hoc accessibility tree only, and never the ground
truth. The restriction is enforced at the data boundary.

Conditions are made physical in a real browser rather than simulated. Each task is
rendered as an HTML form whose submit handler writes the truly committed values to
an authoritative store and, separately, renders a confirmation region into the
accessibility tree. Under deceptive injection the handler commits the attacker's
value while the confirmation exposes the benign one. A pluggable agent perceives the
page through its accessible representation, decides the field values, performs them,
and writes its own narration. The rig then reads ground truth from the authoritative
store and reads the accessible view the way a screen reader would, by resolving each
field through its accessible name and by capturing the full accessibility snapshot.
The agent supplies the narration but never touches the committed state.

Figure~\ref{fig:diverge} shows a captured example. Under deceptive injection the
committed recipient is a homograph of the benign address (digit one for the letter
l) while the accessibility snapshot exposes the benign address. The two channels
genuinely disagree, and the disagreement is read from a production accessibility
snapshot, establishing that the threat model is realisable rather than assumed.

\begin{figure}[t]
\centering
\fbox{\parbox{0.86\linewidth}{\small\ttfamily
=== deceptive\_injection (live browser) ===\\
ground truth : email.recipient = accounts@northgate-1egal.co.uk\\
a11y tree ~~~: email.recipient = accounts@northgate-legal.co.uk\\
channels disagree? YES
}}
\caption{A single live run. The committed recipient is hijacked by a homograph; the
accessibility tree, read from a real snapshot, still speaks the benign address.}
\label{fig:diverge}
\end{figure}

\subsection{Cross-Checking Against a Screen Reader}
The rig reads the accessible channel through a browser accessibility snapshot, a
faithful model of the accessibility tree but not identical to what a production
screen reader announces. To validate the model, Baseerat includes a cross-check that
parses an NVDA Speech Viewer log into the same field-value representation and
compares it against the snapshot channel. The parser and comparator are released and
tested against captured logs; the live capture is performed on a Windows virtual
machine following a documented procedure. Under deceptive injection the two
accessible surfaces agree with each other on the benign spoofed value and both
diverge from ground truth, which is precisely why the accessible channel is
insufficient.

\section{Metrics}

\paragraph{Self-report fidelity.} Under a claim of success every field is implicitly
claimed correct, so fidelity is the fraction of expected fields whose true committed
value really is correct in ground truth. It is computed from ground truth and is
independent of any auditor.

\paragraph{Non-visual detectability.} Given only the non-visual channel, the
\emph{recall} over truly-diverged runs (the fraction that should have been caught
that were), and separately the \emph{false-alarm rate} on clean runs.

\paragraph{Oversight parity gap.} The headline metric. Let $D_{\mathrm{vis}}$ and
$D_{\mathrm{nv}}$ be the detection rates on truly-diverged runs under the visual and
non-visual channels respectively, over the same runs. Then
$\mathrm{Gap} = D_{\mathrm{vis}} - D_{\mathrm{nv}}$. A positive gap is oversight that
a sighted overseer can exercise and a blind one structurally cannot.

\paragraph{Defence effect.} With a trusted receipt present the non-visual detection
rate rises to $D_{\mathrm{nv}}^{\mathrm{rcpt}}$ and the residual gap is
$D_{\mathrm{vis}} - D_{\mathrm{nv}}^{\mathrm{rcpt}}$; the gap closed by the defence
is the difference between the narration-only gap and the residual gap.

\section{Reference-Auditor Results}
\label{sec:results}

We first report a deterministic \emph{reference auditor} that encodes the single
most obvious cross-check a diligent overseer would perform: does the accessibility
tree corroborate the agent's claim of success? On the visual channel it checks the
claim against ground truth. The reference auditor is not intended to be intelligent;
it establishes the benchmark's construct validity and serves as the scalable lower
bound against which the language-model and human auditors are compared. Its results,
over the full 150-task corpus (450 runs), are exact and reproducible.

\begin{table}[t]
\caption{Reference-auditor results over 150 tasks (450 runs). The parity gap is 0.50
overall and concentrated entirely in the deceptive-injection condition. No false
alarms on clean runs.}
\label{tab:reference}
\begin{tabular}{@{}lccc@{}}
\toprule
\textbf{Quantity} & \textbf{clean} & \textbf{silent failure} & \textbf{deceptive
injection} \\
\midrule
Self-report fidelity & 1.00 & 0.61 & 0.61 \\
Visual detection rate & n/a & 1.00 & 1.00 \\
Non-visual detection rate & n/a & 1.00 & 0.00 \\
\textbf{Oversight parity gap} & 0.00 & 0.00 & \textbf{1.00} \\
\bottomrule
\end{tabular}
\end{table}

Table~\ref{tab:reference} shows the result. The non-visual channel catches every
silent failure and no deceptive injection, exactly as the condition design predicts:
when the accessibility tree is honest the cross-check succeeds, and when it is
poisoned in step with the narration the cross-check is defeated. The overall gap is
0.50 and lives entirely in the deceptive-injection condition. Because the poisoning
defeats the reference auditor regardless of how the hijacked value differs from the
intended one, the reference auditor is archetype-invariant; the archetypes are
provided so that the language-model auditor can be analysed by subtlety.

\section{The Defence: Trusted Action Receipts}

A \emph{trusted action receipt} is a machine-verified, accessibly-rendered statement
of what was actually committed, produced by an instrumented layer that observes the
committed state directly and is independent of the page's own confirmation. In the
rig the receipt is rendered from the authoritative committed store, so under
deceptive injection the receipt reports the true hijacked value while the
confirmation still lies. An auditor given a receipt treats it as authoritative and
cross-checks it against the expected values.

\begin{table}[t]
\caption{The defence over 150 tasks. A trusted receipt channel restores non-visual
detection from 0.50 to 1.00 and closes the parity gap to zero, without false alarms
on clean runs.}
\label{tab:defence}
\begin{tabular}{@{}lcc@{}}
\toprule
& \textbf{narration only} & \textbf{with receipt} \\
\midrule
Non-visual detection rate (diverged runs) & 0.50 & \textbf{1.00} \\
Oversight parity gap & 0.50 & \textbf{0.00} \\
\bottomrule
\end{tabular}
\end{table}

Table~\ref{tab:defence} shows that the receipt restores non-visual detection to 1.00
and closes the gap to zero, with no false alarms on clean runs, and the same result
holds when the runs are produced by the live browser rig rather than the
deterministic simulator. A trusted channel, independent of the page and rendered
into the accessible modality, restores to a blind overseer exactly the detection a
sighted overseer had. The realism caveat is explicit: the defence assumes such an
instrumentation layer exists and cannot itself be spoofed by the page. Baseerat does
not claim that assumption is free; it measures how much oversight the assumption
would buy, which is the whole of the gap.

\section{Protocol for the Auditor Arms}

\paragraph{Language-model auditor sweep.} A language model is placed under the same
perceptual restriction as each channel and asked for a structured verdict. This is
not a simulation of a blind person; it is a model of the oversight \emph{channel},
the information bottleneck the gap measures. The released rig runs this arm across
models; the corpus stratification supports a per-archetype and per-domain breakdown
of where the non-visual channel is most and least detectable. Table~\ref{tab:llm}
gives the reporting template.

\begin{table}[t]
\caption{Reporting template for the language-model auditor sweep. The reference row
is the deterministic lower bound reported here; the model rows are populated by the
released sweep.}
\label{tab:llm}
\begin{tabular}{@{}lcccc@{}}
\toprule
\textbf{Auditor model} & \textbf{Vis. recall} & \textbf{Non-vis. recall} &
\textbf{Parity gap} & \textbf{Gap w/ receipt} \\
\midrule
reference (deterministic) & 1.00 & 0.50 & 0.50 & 0.00 \\
% TABLE3-ROWS-START (auto-filled by paper/fill_table3.py; do not edit by hand)
language model A & pending & pending & pending & pending \\
language model B & pending & pending & pending & pending \\
language model C & pending & pending & pending & pending \\
% TABLE3-ROWS-END
\bottomrule
\end{tabular}
\end{table}

\paragraph{Human expert-auditor arm.} Because a language model under a perceptual
restriction is a lower bound and not a person, the language-model arm is anchored by
an expert-auditor evaluation with consenting professional blind accessibility
auditors, who audit a stratified sample of runs through the non-visual channel. This
arm calibrates the language-model auditor against human non-visual detection and is
the basis for the present work's central empirical claim.

\section{Limitations and Threats to Validity}

\paragraph{The agent is scripted for the reference results.} In the reference
configuration the agent performs the task correctly and the divergence is introduced
by the environment. This is a deliberate control that isolates the oversight
question from agent competence, so that a measured gap reflects the channel and not
agent error. A real computer-use agent slots into the same seam.

\paragraph{The receipt defence assumes a trusted layer.} The defence's power depends
on an instrumentation layer that the page cannot spoof. Where no such layer exists,
the measured gap stands as the exposure.

\paragraph{The accessible channel is modelled by a snapshot.} The rig reads the
accessibility tree through a browser snapshot. The NVDA cross-check validates this
against a production screen reader on a stratified sample; a full screen-reader in
the loop for every run is future work.

\section{Ethics and Accessibility Statement}

This work is motivated by a safety exposure that falls disproportionately on blind
and low-vision users, and is intended to reduce it. The corpus contains only
fictional values and no real personal data. The deceptive-injection condition is a
controlled construction inside a local rig and provides no attack against any
deployed system. The human expert-auditor arm is designed as a paid, non-clinical,
minimal-risk evaluation with informed consent.

\section{Conclusion}

Agent oversight has a hidden precondition: that the overseer can see the screen. For
a blind screen-reader user that precondition fails, and the oversight channel
collapses onto an agent self-report that the literature has shown to be unreliable,
and that an adversary can actively poison in step with the action. Baseerat turns
this into a measurement. Over a 150-task corpus a deterministic reference auditor
exhibits an oversight parity gap of 0.50, concentrated entirely in a
deceptive-injection condition that we demonstrate is realisable in a live browser,
and a trusted-receipt defence closes the gap to zero. Model developers can and
should be scored on the oversight they afford an overseer who cannot see the screen,
not on the oversight they afford one who can.

\begin{acks}
The benchmark, the live rig, the task generator, the corpus, the defence, the NVDA
cross-check, and the reference results are released as an open repository. This work
is archived at doi:10.5281/zenodo.22080740. Copyright WAJD Group. Developed by WAJD
AI.
\end{acks}

\begin{thebibliography}{9}
\bibitem{osharm} OS-Harm: A Benchmark for Measuring Safety of Computer Use Agents.
arXiv:2506.14866.
\bibitem{silentfailure} From Confident Closing to Silent Failure: Characterizing
False Success in LLM Agents. arXiv:2606.09863.
\bibitem{mirage} MIRAGE-Bench: LLM Agent is Hallucinating and Where to Find Them.
arXiv:2507.21017.
\bibitem{a11ycua} A11y-CUA: Characterizing the Accessibility Gap in Computer Use
Agents. arXiv:2602.09310.
\bibitem{xaiblv} Explainable AI for Blind and Low-Vision Users: Navigating Trust,
Modality, and Interpretability in the Agentic Era. arXiv:2604.00187.
\bibitem{nexui} Navigation Alone Is Not Enough: Evaluating Explanatory Assistive UI
Agents (NeXUI). arXiv:2608.09944.
\bibitem{morae} Morae: Proactively Pausing UI Agents for User Choices.
arXiv:2508.21456.
\bibitem{receipts} Tool Receipts, Not Zero-Knowledge Proofs: Practical Hallucination
Detection for AI Agents. arXiv:2603.10060.
\bibitem{misscore} How Benchmarks Mis-Score Computer-Use Agents. arXiv:2607.28367.
\end{thebibliography}

\end{document}
